% $MirOS: contrib/code/Flyers/en.tex,v 1.3 2012/01/27 15:39:20 bsiegert Exp $
%

\documentclass[a4paper,landscape,11pt,notumble]{leaflet}
\usepackage{mirosflyer}

\hyphenation{Mir-OS}

\begin{document}

%\maketitle
%\thispagestyle{empty}

\graybox
\begin{center}
{\noindent\raisebox{-1ex}{\includegraphics[height=3ex]{img/Flag_of_the_United_Kingdom}}%
\quad E~~N~~G~~L~~I~~S~~H\par}
\vfill

\fontsize{64}{64}\selectfont WTF\par
\fontsize{30}{30}\selectfont is the\par
\fontsize{56}{56}\selectfont MirOS\par
Project?\par

\vfill

\includegraphics[width=0.75\columnwidth]{img/m-color}

\vfill%
\end{center}%
\graybox

\newpage

\section{Current Version}

\subsection{MirOS \#10-current}

The CD distributed at this event contains MirOS \#10-current---a beta version of MirOS \#11. This system is based on OpenBSD-current and older versions. MirOS BSD also includes an improved shell (mksh), ports tree (MirPorts), acronym database (wtf) and GNU cvs 1.12 with custom extensions.

The \textcolor{darkred}{bootloader} has been mostly rewritten and is now multiboot compliant so it can be called from GRUB.
The \textcolor{darkred}{random number generator} (SRNG) has been improved, and many core components have been updated.
pkgsrc® can be installed in parallel to and as a complement of MirPorts, giving access to thousands of packages for current software. 

The included gcc 3.4 compiler has the Propolice Stack Smashing Protector and supports C, C++, Pascal,  Objective-C and Ada.


\subsection{About the DVD}

The DVD distributed at this event allows installing MirOS~BSD on the \textcolor{darkred}{i386 and sparc architecture}.
It does not contain a live system; however, live CDs can be downloaded from our homepage.
To start the installation, boot from the DVD and select ``Install''.

The DVD also contains a repository of \textcolor{darkred}{over 5000 binary packages} (for i386) compiled from a stable pkgsrc® branch.
To use these packages, select the \url{pkgsrc10.ngz} set during the installation.
Then, add \url{/usr/pkg/bin} and \url{/usr/pkg/sbin} to your PATH and execute:
%
\begin{verbatim}
export PKG_PATH=/cdrom/packages/All
\end{verbatim}
%
You can now install binary packages with a simple pkg\_add command, for example ``sudo pkg\_add screen''.

\newpage

\section{MirOS BSD}

\largered{MirOS BSD is a secure computer operating system from the BSD family.}

It is a derivative of OpenBSD. Source code from OpenBSD is regularly imported and merged. A lot of code and ideas are adopted from NetBSD® and other sources.

MirOS was started after some differences in opinion between Theo de~Raadt, the OpenBSD project leader, and Thorsten Glaser, who is now our lead developer. The principal maintainer of MirPorts is Benny Siegert. Several more persons are contributors for the project.

MirOS BSD often anticipates bigger changes in OpenBSD and includes them before OpenBSD itself. For example, ELF on i386 and support for gcc3 were available in MirOS first. Controversial decisions are often made differently from OpenBSD; for instance, SMP as implemented in OpenBSD is not supported in MirOS.

Most of the differences between OpenBSD and MirOS are numerous small details. The end result is, hopefully, a more refined BSD experience.

The base system has been trimmed down. Seldom used components like NIS, Kerberos, Bind and the BSD games have been removed. The last two are installable as ports. Support for internationalisation, Unicode and UTF-8 is integrated, Citrus libiconv is also available.

For the stable releases, binary security updates are available and can be directly selected during install.

IPv6 is supported everywhere, even in the Apache 1.3-based web server. ISDN drivers are also included.

We are legally bound to point out the advertising clauses. They are too many to print them here but you can find them at \url{http://www.mirbsd.org/about.htm}.

\newpage

%%% Back Cover Page

\graybox\par
\begin{center}
\vspace{3ex}
{\Large\itshape
Enjoy MirOS!\par\medskip
The developer team\par\medskip
\color{darkred}%
Thorsten Glaser\\
Benny Siegert\par
}

\vfill

\includegraphics[width=0.75\columnwidth]{img/bsdaemon}

\vfill%
\end{center}%
\subsection{Contact}
{\noindent\large
\begin{tabular}{ll}
home page: & \textcolor{darkred}{\url{http://www.mirbsd.org/}}\\
%& http://www.mirbsd.org/\\
IRC: & \textcolor{darkred}{\url{#mirbsd}} or \textcolor{darkred}{\url{#!/bin/mksh}}\\
& on irc.freenode.net\\
\end{tabular}
\par}
\vspace{3ex}%
\graybox

\newpage

\section{mksh, the MirOS Korn Shell}

The standard shell in MirOS is mksh, derived from the Korn Shell. bash users should immediately feel at home as the commands are very similar.

mksh is smaller and much faster than other shells, such as bash or zsh. It contains the features from the OpenBSD and Debian pdksh and many, many bugfixes. Its code is very clean (warnings, const cleanup) and has been checked for security problems by the developers and by Coverity, Inc.

mksh is great for interactive use but lends itself equally well for shell scripting purposes. Even with its small size, no important features are missing. UTF-8 is fully supported, of course.

mksh is portable and runs on a wide variety of architectures and operating systems, e.\,g. BSD, Linux, Mac OS, Solaris, HURD, HP-UX, AIX, IRIX, QNX, Minix 3, Tru64, Ultrix and even Windows. The following systems (and more) come with an mksh package:

\begin{itemize}
\item Debian, Gentoo, FreeWRT, Fedora and other GNU/Linux distributions, OpenSUSE Buildservice
\item NetBSD® pkgsrc®, FreeBSD® Ports, MidnightBSD
\item fink, MacPorts and other Mac toolkits
\end{itemize}

\section{Installing additional software}

For installing additional software not in the base system, we offer two solutions: 
\textcolor{darkred}{MirPorts}---a derivative of the OpenBSD ports tree---and \textcolor{darkred}{pkgsrc®}, which is maintained by the NetBSD project.
MirOS~BSD support for pkgsrc® is a work in progress; however, it contains more packages (more than 11000) and in newer versions.
Binary packages for pkgsrc® are available on \url{ftp.netbsd.org}.

\newpage

\section{Introduction to MirPorts}

After the first checkout or after updates, ``make setup'' in \url{/usr/ports} installs the package tools and configures the MirPorts infrastructure. The ports themselves are in subdirectories, sorted by category. Just executing ``mmake install'' in such a directory will download the source code, compile it, create a binary package and install it. Dependencies are automatically installed when necessary. Some ports exist in different ``flavours'', e.\,g. with or without X support.

Both MirOS and MirPorts should put most of the ``dotfiles'' in users' home directories in a single directory named ``.etc''. You can have your own programs and scripts in \url{~/.etc/bin}.

MirPorts is portable. It has support for the following operating systems:

\begin{itemize}
\item MirOS BSD (-stable and -current)
\item OpenBSD (-stable and -current)
\item MidnightBSD
\item Mac OS X (10.4 and newer) / Darwin
\item Interix / SFU 3.5 (limited support)
\end{itemize}

\noindent
Even on stable releases, using the newest MirPorts version is recommended. 

\subsection{MirLibtool}

Libtool is used by many packages to build shared libraries in  a portable way. However, there are many problems with it—for example, it breaks when no C++ compiler is installed. Therefore, MirPorts contains a modified version named MirLibtool.

MirLibtool is based on GNU libtool 1.5. It is compatible with all versions of autotools. The MirPorts infrastructure installs it automatically whenever a port uses autoconf to recreate its configure script.

\end{document}
